% ============================================================
%  Interpreting FCFS Appointment Simulation Results
%  Under the Current Formulation
%  Metric behavior, access losses, and model checks
% ============================================================

\documentclass[11pt]{article}

\usepackage[margin=0.9in]{geometry}
\usepackage{amsmath}
\usepackage{amssymb}
\usepackage{booktabs}
\usepackage{graphicx}
\usepackage{float}
\usepackage[hidelinks]{hyperref}
\usepackage{array}
\usepackage[table]{xcolor}
\usepackage{microtype}
\usepackage{parskip}
\usepackage{enumitem}
\usepackage{mdframed}
\usepackage{tabularx}

\graphicspath{%
  {../../figures/}%
  {../../../../../outputs/class1_balking/figures/}%
  {../../../../../outputs/class1_balking_threshold/figures/}%
}

\sloppy
\setcounter{tocdepth}{2}

\newcommand{\mainfigure}[3]{%
\begin{figure}[H]
    \centering
    \includegraphics[width=\textwidth]{#1}
    \caption{#2}
    \label{#3}
\end{figure}%
}

\newmdenv[backgroundcolor=yellow!12, linecolor=gray!60, linewidth=0.6pt,
          innertopmargin=6pt, innerbottommargin=6pt]{checkenv}

\title{Interpreting FCFS Appointment Simulation Results\\
Under the Current Formulation\\[0.4em]
\large Metric behavior, access losses, and model checks}
\author{CUIMC Appointment Simulation}
\date{May 2026}

\begin{document}
\maketitle

\begin{abstract}
\noindent
This note interprets simulation outputs from the first discrete-event FCFS appointment
model. The goal is to understand what the outputs show and whether they are
mechanically plausible under the current formulation.


The central finding is that utilization and served rate separate sharply.
At baseline (\textit{$\lambda = 100$ arrivals per day, 32 slots per day, 14-day
horizon}), average utilization is 0.839 while overall served rate is 0.269.

This separation is mechanically plausible because the two metrics have
different denominators: utilization is normalized by slots, served rate by
arrivals. The baseline is structurally oversubscribed. Patients are indeed served
on future days (the horizon extends to $H = 14$ days), but each service day
has exactly 32 slots regardless of which arrival cohort fills them. Over
any measurement window of $D$ days, at most $32D$ visits can be completed
against approximately $100D$ arrivals, so served rate cannot exceed
$32/100 = 0.32$.

Arrival outcome decomposition at baseline shows that balking (35.6\% of
arrivals) and cancellation (32.3\%) are the dominant loss channels, not
horizon saturation: no-offer is~0 at baseline demand. The high cancellation
fraction is mechanically consistent with the long average booking delay
(8.35~days) combined with a 10\%~daily cancellation probability.
A controlled no-balking diagnostic confirms that suppressing balking
increases booking conversion but not completed visits: unmet demand shifts
toward no-offer and cancellations, while completed utilization remains
pinned near~0.84.

Mean offered wait is preferable to mean accepted wait as a delay measure,
because offered delay is recorded before the balking decision and therefore
includes patients who ultimately reject the offer. Offered wait still excludes
no-offer patients, so it must be read together with the no-offer share; which
is zero at baseline but grows above $\lambda \approx 110$.
\end{abstract}

\tableofcontents
\newpage

%% ============================================================
\section{Simulation Mechanics Relevant for Interpretation}
\label{sec:mechanics}
%% ============================================================

The following mechanics directly affect how outputs should be read. The full
formulation remain in the simulation documentation.

\paragraph{Finite booking horizon and FCFS pooling.}
Each day, arrivals for both classes are pooled and randomly permuted into one
FCFS sequence. Each patient is offered the earliest available slot within the
$H$-day horizon (including same-day, $r=0$). Because FCFS operates on a
permuted joint sequence, class labels are irrelevant to slot assignment:
neither class is prioritized. Class disparities in outcomes arise only from
differences in arrival rates or behavioral parameters.

\paragraph{Offered delay $\tau$ recorded before balking.}
When a valid slot exists, the offered delay $\tau$ is recorded before the
patient decides whether to accept. Mean offered wait therefore includes
patients who ultimately balk. Mean accepted wait excludes them, which biases
it downward when high-$\tau$ patients disproportionately balk.

\paragraph{Future cancellations only; no rebooking of no-show slots.}
Cancellations apply only to future appointments ($r \geq 1$). Same-day
cancellations are not modeled. Every future booking receives a fresh
Bernoulli draw with probability \texttt{cancel\_prob} on each day before
service; at mean accepted wait $8.33$~days and \texttt{cancel\_prob}~$=0.10$,
survival to service is approximately $0.9^{8.33}=0.416$. Slots vacated by
no-shows are not rebooked---once a patient no-shows, that slot is lost.

\paragraph{No-offer rule.}
If the horizon is full when a patient reaches the front of the FCFS queue,
that patient and all remaining arrivals for that day are counted as
\texttt{no\_offer}. This means a single saturation event affects the entire
remainder of the day's queue, which could introduce daily-order artifacts.

\paragraph{Utilization denominator.}
Utilization counts completed visits per available slot, not scheduled
appointments per slot. No-shows, cancellations, and unbooked slots all reduce
utilization. A high-utilization day is one where most slots became completed
visits.

%% ============================================================
\section{Baseline Setup and Diagnosis}
\label{sec:baseline}
%% ============================================================

\subsection{Parameter Setup}

The baseline parameters are listed in Figure~\ref{fig:baseline-params-and-step-functions}. The key feature is the demand-to-capacity ratio: \(100\) arrivals per day against \(32\) slots per day, or roughly a \(3{:}1\) oversubscription. This is not meant to represent a plausible steady-state clinic. It is a stress-test baseline, chosen to make the booking queue fill and to expose how the FCFS simulation metrics behave under constrained capacity.

At baseline, the two classes are intentionally symmetric: they have the same arrival rate, balking rule, cancellation probability, and no-show rule. This symmetry makes the baseline a neutral reference point. The sensitivity analysis then perturbs one class-specific parameter at a time to study both the direct effect on the modified class and the spillover effect on the other class through the shared FCFS booking system. The goal is to interpret the mechanics of the current formulation, not to claim that the baseline is empirically calibrated.

\begin{figure}[H]
\centering

\begin{minipage}[t]{0.4\textwidth}
\vspace{0pt}
\renewcommand{\arraystretch}{1.15}
\small
\begin{tabularx}{\linewidth}{Xr}
\toprule
Parameter & Value \\
\midrule
Class 1 arrival rate $\lambda_1$ (per day) & 50 \\
Class 2 arrival rate $\lambda_2$ (per day) & 50 \\
Total arrival rate $\lambda$ (per day) & 100 \\
Slots per day $S$ & 32 \\
Booking horizon $H$ (days) & 14 \\
Cancellation probability $\phi$ & 0.10 \\
Balking threshold $\theta^{\text{balk}}$ (days) & 9 \\
High-delay balking prob.\ $b^{\text{high}}$ & 0.50 \\
Low-delay balking prob.\ $b^{\text{low}}$ & 0.00 \\
No-show threshold $\theta^{\text{no-show}}$ (days) & 6 \\
High-delay no-show prob.\ $\xi^{\text{high}}$ & 0.30 \\
Low-delay no-show prob.\ $\xi^{\text{low}}$ & 0.00 \\
\bottomrule
\end{tabularx}
\end{minipage}%
\hfill
\begin{minipage}[t]{0.58\textwidth}
\vspace{0pt}
\centering
\includegraphics[width=\linewidth,height=0.82\linewidth,keepaspectratio]{baseline_step_functions.png}
\end{minipage}

\caption{Left: baseline parameter values (\texttt{configs/baseline.yaml}),
symmetric across both classes; measurement window 365 days after a 30-day
burn-in. Right: balking and no-show step functions at baseline. Both
probabilities are~0 below their respective thresholds and jump to 0.50
(balking) and 0.30 (no-show) above them.}
\label{fig:baseline-params-and-step-functions}
\end{figure}

\subsection{Aggregate Metrics at Baseline}

Metrics are averaged over a single run of 365 measurement days,
preceded by a 30-day burn-in to reach steady state and followed by a
10-day cooldown to allow late bookings to resolve. All figures in
Table~\ref{tab:baseline} are point estimates from that single run; the
sensitivity appendix uses multiple randomization seeds to assess stability and the influence of each parameter.

\begin{table}[H]
\centering
\renewcommand{\arraystretch}{1.15}
\begin{tabular}{lr}
\toprule
Metric & Value \\
\midrule
Average utilization & 0.839 \\
Overall served rate & 0.269 \\
Mean offered wait (days) & 9.30 \\
Mean accepted wait (days) & 8.35 \\
Overall balking rate & 0.356 \\
Class served-rate gap ($\Delta_{\text{access}}$) & 0.001 \\
\bottomrule
\end{tabular}
\caption{Baseline aggregate metrics. The near-zero class gap is expected:
both classes have identical behavioral parameters and symmetric arrival rates,
so class labels are exchangeable under pooled FCFS.}
\label{tab:baseline}
\end{table}

\subsection{The Utilization--Access Gap}

\[
  \text{average utilization}
  = \frac{\text{completed visits}}{\text{available slots}},
  \qquad
  \text{overall served rate}
  = \frac{\text{completed visits}}{\text{arrivals}}.
\]

Utilization~(0.839) and served rate~(0.269) differ by a factor of roughly
3.1. This is mechanically expected because the denominators differ.
Utilization is normalized by supply; served rate is normalized by demand.
When demand far exceeds supply, served rate is bounded by the
supply-to-demand ratio; utilization is not.

At baseline, $100/32 = 3.125$~arrivals per slot. The ceiling argument is
a flow constraint: patients are served on future days (via the booking
horizon), but each service day has exactly 32 slots regardless of which
arrival cohort fills them. Over $D$ days at most $32D$ visits can complete
against $100D$ arrivals, so served rate cannot exceed $1/3.125 = 0.32$.
The observed served rate of 0.269 falls below this ceiling because
no-shows consume slots without producing completed visits.

The simulation implies, under these assumptions, that a facility-level
utilization measure would give a misleading picture of patient access at
this demand-to-capacity ratio.

%% ============================================================
\section{Arrival Outcome Decomposition}
\label{sec:decomp}
%% ============================================================

\subsection{Loss Channels at Baseline}

Every arrival ends up in exactly one outcome state. At baseline:

\begin{table}[H]
\centering
\renewcommand{\arraystretch}{1.12}
\begin{tabular}{lrp{0.52\textwidth}}
\toprule
Outcome & Share of arrivals & Interpretation \\
\midrule
Served & 0.269 & Completed visit \\
Balked & 0.356 & Received an offer, rejected the delay \\
Canceled & 0.323 & Booked, but canceled before service day \\
No-show & 0.052 & Kept booking, did not attend \\
No-offer & 0.000 & Horizon full; no slot offered \\
Unresolved & $<$0.001 & Booked but not resolved before simulation end \\
\midrule
Total & 1.000 & \\
\bottomrule
\end{tabular}
\caption{Arrival outcome decomposition at baseline ($\lambda = 100$).
No-offer is zero: the horizon does not completely fill at baseline demand
because cancellations continuously free future capacity.}
\label{tab:decomp}
\end{table}

Two findings stand out. 
\begin{itemize}
    \item First, the no-offer rate is zero even at 100
arrivals/day. The horizon does not saturate because cancellations
(10\%/day for each future appointment) continuously free slots. However,
cancellation frees far-future capacity, which is then offered to new
arrivals at long delays---reinforcing the same long-delay dynamic that
drives balking.

\item Second, the 32.3\% cancellation rate (as a share of arrivals) is not
self-evidently wrong. Of the 64.4\% of arrivals that book an appointment,
roughly half ultimately cancel. With an average booking delay of 8.35~days
and a 10\%~daily cancellation probability, the implied cumulative
cancellation probability is $1 - 0.9^{8.35} \approx 0.59$. The observed
cancellation rate among booked patients is in the same range, though
somewhat lower; possibly because some booked patients face fewer
cancellation-check days when scheduled soon.
\end{itemize}

The simulation engine applies a
fresh Bernoulli draw to every future appointment on each simulated day, so the
cumulative argument is consistent with the implementation.

\subsection{Balking Suppression Diagnostic}
\label{sec:balk-suppression}

A controlled 30-seed rerun shows that suppressing balking does not increase
completed utilization. Average utilization changes only from \(0.8418\) to
\(0.8404\), with a 95\% confidence interval for the difference of
\([-0.0029,0.0001]\). Overall served rate decreases slightly, from
\(0.2692\) to \(0.2682\).

The mechanism is a channel shift rather than a utilization gain. Removing
balking increases the booked rate by \(5.85\) percentage points, but the
booking horizon fills faster, generating a no-offer share of \(0.2966\). It
also increases cancellation exposure because more patients accept long-delay
bookings, raising the cancellation share from \(0.3248\) to \(0.3836\).
Therefore, additional booking conversion is almost exactly offset by no-offer
accumulation and cancellations. Under the current formulation, balking does
not materially raise or lower utilization; it determines whether unmet demand
appears as balking or no-offer.

\begin{table}[H]
\centering
\small
\renewcommand{\arraystretch}{1.08}
\begin{tabularx}{\textwidth}{>{\raggedright\arraybackslash}Xcccc}
\toprule
Metric & Baseline & No-balking & Difference & 95\% CI \\
\midrule
Average utilization & 0.8418 & 0.8404 & -0.0014 & [-0.0029, +0.0001] \\
Overall served rate & 0.2692 & 0.2682 & -0.0011 & [-0.0018, -0.0003] \\
Booked rate & 0.6450 & 0.7034 & +0.0585 & [+0.0567, +0.0602] \\
No-show share & 0.0506 & 0.0507 & +0.0002 & [-0.0003, +0.0007] \\
Balked share & 0.3550 & 0.0000 & -0.3550 & -- \\
Cancelled share & 0.3248 & 0.3836 & +0.0588 & -- \\
No-offer share & 0.0000 & 0.2966 & +0.2966 & -- \\
Mean accepted wait & 8.33 d & 9.45 d & +1.12 d & -- \\
\bottomrule
\end{tabularx}
\caption{Controlled balking-suppression diagnostic over 30 seeds.
Suppressing balking increases booking conversion but does not increase
completed utilization; additional bookings are offset by no-offer
accumulation and higher cancellation exposure.}
\label{tab:balk-suppression}
\end{table}

\mainfigure{../../../../../figures/balking_effect_verification.png}{%
Controlled sweep of high-delay balking probability with all other parameters
held at baseline. The baseline marker is at high-delay balking probability
\(0.5\); moving to zero removes balking, but average utilization and overall
served rate remain nearly flat within simulation uncertainty.
}{fig:balk-suppression-sweep}

\subsection{How Loss Channels Shift with Demand}

Figure~\ref{fig:decomp-full} shows the arrival outcome decomposition as total
daily arrivals vary from 30 to 170. At low demand ($\lambda \lesssim 50$),
virtually all arrivals are served. As demand rises above the 32-slot daily
capacity, the served rate falls. The no-offer rate remains zero up to
$\lambda \approx 110$, at which point the horizon begins to saturate.

The transition from a cancellation- and balking-dominated regime (below
$\lambda \approx 110$) to a horizon-saturation regime (above) is a
structural feature of the model. It suggests that reporting the no-offer
share alongside served rate and offered wait is necessary to understand
which failure mode dominates in a given parameter setting.

\mainfigure{fcfs_arrival_outcome_decomposition.png}{%
Arrival outcome decomposition as total daily arrivals vary from 30 to 170.
At baseline ($\lambda = 100$, marked by a dashed vertical line), balking and
cancellation are the dominant loss channels; no-offer is zero. No-offer
becomes substantial only above $\lambda \approx 110$.}{fig:decomp-full}

%% ============================================================
\section{Cross-Metric Diagnostics}
\label{sec:crossmetric}
%% ============================================================

Figure~\ref{fig:main-panel} shows three diagnostics that together
capture the core structure of the simulation output.

\mainfigure{research_note_main_panel.png}{%
\textit{Left:} Arrival outcome decomposition as total daily arrivals
vary. At baseline ($\lambda = 100$, dashed line), balking and
cancellation are the dominant loss channels; no-offer is
zero. No-offer appears only above $\lambda \approx 110$.
\textit{Center:} Utilization--access frontier. The plot reveals two
regimes along the demand sweep. At higher demand ($\lambda \geq 60$, the
main cluster), utilization is pinned near~0.84 while served rate falls
from~0.45 to~0.16. At lower demand ($\lambda \leq 50$, the three isolated
points), both metrics co-move.
\textit{Right:} Offered vs.\ accepted wait. Points below the~45$^\circ$
line indicate that accepted wait understates the delay patients face.
The gap widens with balking rate (red intensity).}{fig:main-panel}

\textbf{1.~Utilization--access frontier (center panel).}
The center panel shows two distinct regimes along the demand sweep. At higher
demand ($\lambda \geq 60$), utilization is pinned near~0.84 while served rate
collapses from~0.45 to~0.16. This should be read as oversubscription plus
daily cancellation and rebooking dynamics, not as evidence that balking
improves capacity use. The controlled no-balking diagnostic in
Section~\ref{sec:balk-suppression} shows that suppressing balking at baseline
does not increase completed utilization; it shifts losses toward no-offer and
cancellations.

\textbf{2.~Offered vs.\ accepted wait (right panel).}
All points fall below the 45$^\circ$ line: accepted wait consistently
understates offered wait. The gap widens as balking rate increases (red
intensity). This is mechanically consistent: patients who balk tend to have
above-average $\tau$, so removing them from the accepted sample pulls the
mean down. Any policy or parameter change that raises balking will make
accepted wait appear to improve while offered wait and served rate both
worsen.

\textbf{3.~No-offer and the boundary of offered-wait validity.}
At baseline, no-offer is zero, so offered wait covers all patients who
received any delay information. Above $\lambda \approx 110$, a growing
share of arrivals receive no offer and are excluded from offered wait
entirely. At $\lambda = 170$, the no-offer share reaches approximately
34.6\% of arrivals. In that regime, offered wait alone misses the
worst-off patients.

%% ============================================================
\section{Metric Hierarchy Under This Formulation}
\label{sec:metrics}
%% ============================================================

\begin{table}[H]
\centering
\renewcommand{\arraystretch}{1.2}
\begin{tabular}{p{0.20\textwidth}p{0.22\textwidth}p{0.24\textwidth}p{0.24\textwidth}}
\toprule
Metric & What it measures & Mechanical interpretation & Main caveat \\
\midrule
Overall served rate & Arrivals $\to$ completed visits & Primary access measure; bounded above by $S/\lambda$ & Does not decompose loss channels \\
Average utilization & Completed visits per available slot & Capacity-use measure; can remain high under oversubscription even while served rate falls & High value does not imply good access \\
Mean offered wait & Delay among patients who received any offer & Less selective than accepted wait; includes patients who then balk & Excludes no-offer patients entirely \\
Mean accepted wait & Delay among patients who accepted bookings & Conditional; pulled down when high-delay patients balk & Can fall as access worsens \\
Balking rate & Share of offered patients who reject & Congestion diagnostic & Not a success metric \\
Class served-rate gap & Class 1 served rate minus Class 2 & Meaningful only when classes differ in behavior or arrivals & Near-zero in symmetric baseline \\
\bottomrule
\end{tabular}
\caption{Metric interpretation under the current FCFS formulation.}
\label{tab:metrics}
\end{table}

\subsection{High-Demand Decoupling and Balking Selection}

The center panel of Figure~\ref{fig:main-panel} separates low-demand points
from the oversubscribed regime. Every point with $\lambda \geq 60$ has a
positive balking rate (ranging from~4.6\% to~36.1\%) and utilization
constrained to~0.839--0.841; every point with $\lambda \leq 50$ has zero
balking, and both utilization and served rate vary more freely. Because
arrival rate is changing at the same time, this demand sweep should not be
read as a controlled no-balking counterfactual for baseline. The direct
diagnostic in Section~\ref{sec:balk-suppression} shows that suppressing
balking at baseline changes the distribution of access losses, not completed
utilization.

Figure~\ref{fig:no-balk-delay} plots both mean offered and mean accepted
delay against the two behavioral thresholds for $\lambda \leq 110$. In the low-demand points ($\lambda \leq 50$), the two series coincide
because every offer is accepted, no balking is observed because the low-demand ensures offered delays inferior to the balking threshold. Once long-delay offers begin to generate
balking, offered delay continues to rise while accepted delay grows more
slowly: the highest-$\tau$ patients are removed from the accepted sample.
This selection effect explains why accepted wait understates congestion in
high-balking settings, but it is a selection artifact rather than a
completed-visit effect.

\mainfigure{no_balk_delay_thresholds.png}{%
  Mean offered (dashed) and accepted (solid) booking delay from $\lambda = 30$
  to $\lambda = 110$, plotted against the no-show threshold ($6$~days, green)
  and balking threshold ($9$~days, purple). Shading marks the no-balking
  ($\lambda \leq 50$, blue) and balking ($\lambda \geq 60$, purple) regimes.
  The two series coincide where no one balks and diverge thereafter, with the
  gap widening as high-delay patients leave the accepted sample.
}{fig:no-balk-delay}

\subsection{Offered Wait vs.\ Accepted Wait}

The right panel of Figure~\ref{fig:main-panel} shows that accepted wait
is consistently below offered wait across all demand levels, with the gap
widening as balking rate rises; Figure~\ref{fig:no-balk-delay} shows the
same split in absolute terms as a function of~$\lambda$.

The implication is that accepted wait understates congestion in
high-balking regimes. Offered wait is the better denominator for any
congestion diagnostic, but it still excludes no-offer patients, who
receive no delay information at all.

\subsection{No-Offer as a Distinct Access Failure}

Mean offered wait is computed among patients who received a slot offer.
When the horizon saturates ($\lambda \gtrsim 110$ in this formulation),
a growing share of arrivals receive no offer and are entirely absent from
the offered-wait calculation. Figure~\ref{fig:no-balk-delay} covers
$\lambda \leq 110$, the range immediately preceding this saturation
regime; at $\lambda = 110$ offered delay has just crossed the balking
threshold and no-offer patients are beginning to accumulate. A useful
diagnostic is to report no-offer share alongside offered wait. Together
they describe both the level and the completeness of delay information.

%% ============================================================
\section{Main Mechanisms Identified by the Sensitivity Analysis}
\label{sec:sensitivity}
%% ============================================================

The sensitivity runs are the main diagnostic tool in this note. They are not
used to estimate real-world causal effects; they are used to identify which
mechanisms are active under the current FCFS formulation. The one-at-a-time
runs show how the system responds when a single parameter is perturbed around
the baseline, while the randomized regression screen summarizes which
perturbations remain predictive across a broader simulated design space. Full
design details and regression tables are reported in
Appendices~\ref{app:sens-design} and~\ref{app:regression}. Unless stated
otherwise, the coefficients discussed below are statistically significant at
the \(5\%\) level.

\subsection{How to Read the Sensitivity Evidence}

The sensitivity analysis combines two complementary designs. The
one-at-a-time (OAT) slices hold all parameters at their baseline values while
varying a single parameter across a fixed grid. They therefore show local
behavior around the stress-test baseline, not general effects across all
possible configurations. In the class-level OAT slices, one class is perturbed
while the other remains at baseline, so the plots show both the direct effect
on the modified class and the spillover to the other class through the shared
FCFS booking pool.

The randomized regression screen serves a different purpose. It samples
\(240\) parameter settings across broad ranges, averages two simulation seeds
per setting, and fits standardized OLS regressions with HC3 robust standard
errors. The coefficients should be read as average associations over the
sampled parameter space, not as local causal effects at the baseline. An OAT slice can reveal a baseline-specific regime, while
the regression screen identifies which parameters dominate across a wider
range of simulated configurations.

Together, the OAT slices provide shape and mechanism; the regression screen
provides a ranking of dominant drivers.

\subsection{Capacity Pressure and Access Loss}

\paragraph{Demand pressure.}
Total arrival rate has the largest absolute standardized regression
coefficient for overall served rate \((\beta=-0.774)\) and is the dominant
driver of mean offered wait \((\beta=+0.576)\). This is the simplest
mechanism in the simulation: more arrivals compete for the same finite booking
capacity. As demand rises, offers move farther into the booking horizon and a
smaller share of arrivals complete service. This confirms that the served-rate
loss is not a low-utilization phenomenon; it is primarily a high-demand,
fixed-capacity phenomenon.

\subsection{Post-Booking Losses}

\paragraph{No-show threshold.}
The average no-show threshold has the largest coefficient for utilization
\((\beta=+0.515)\). This is mechanically direct. Raising the threshold delays
entry into the high no-show region, so fewer booked slots are lost on the
service day. Since no-show slots are not rebooked, no-shows reduce completed
utilization rather than merely shifting capacity to another patient.

\paragraph{Cancellation and rebooking.}
Cancellation behaves differently from no-show. Cancellation probability has a
positive coefficient for overall served rate \((\beta=+0.116)\) and a
negative coefficient for mean offered wait \((\beta=-0.442)\). This is
counterintuitive only if cancellations are treated as pure losses. In the
implementation, cancellations occur at the start of each day and free future
slots before new arrivals are booked. Under high demand, those freed slots are
likely to be reabsorbed before the service date. Replacement patients often
receive nearer appointments and therefore have less time to cancel or
no-show.

Figure~\ref{fig:cancel-slice} shows the class-level version of this mechanism.
When one class's cancellation probability rises, that class's own served rate
falls because more of its booked patients cancel. The other class's served
rate rises because it can capture some of the freed capacity. The overall
effect is small but positive in the high-load regime. This should not be read
as a general claim that cancellations improve access; it is a consequence of
the current high-demand FCFS setting, where freed future slots are usually
reabsorbed.

\mainfigure{cancellation_probability_slice_access.png}{%
Served rate as a function of one class's cancellation probability, with the
other class held at baseline \(\phi=0.10\). The canceling class loses served
rate, while the other class gains by filling some freed slots. The modest
increase in overall served rate is consistent with a high-load rebooking-chain
effect. Dashed vertical line marks the baseline.
}{fig:cancel-slice}

\subsection{Behavioral Rejection and Loss-Channel Shifts}

\paragraph{Balking and no-show step intensity.}
The average of the high-delay balking and no-show probabilities has negative
coefficients for utilization \((\beta=-0.288)\), overall served rate
\((\beta=-0.106)\), and mean offered wait \((\beta=-0.181)\). The first two
effects are access losses: stronger behavioral frictions reduce the share of
arrivals who complete service. The negative wait coefficient should not be
interpreted as an improvement. It reflects selection: high-delay patients are
more likely to leave through balking or no-offer, so the observed wait metrics
can fall even as access deteriorates.

\paragraph{Balk-vs-no-show dominance.}
The difference between the balking step and the no-show step has positive
coefficients for utilization \((\beta=+0.400)\) and overall served rate
\((\beta=+0.186)\). This does not mean that balking is beneficial in an
absolute sense. It means that, conditional on behavioral intensity, losses
before booking are less damaging to completed slot fill than losses after a
slot has been reserved. If no-show dominates, patients accept long-delay
appointments and then miss them, wasting booked capacity on the service day.
If balking dominates, more losses occur before booking, leaving capacity
available to be offered to other arrivals.

The controlled no-balking diagnostic is consistent with this interpretation.
Suppressing balking entirely increases the booked rate, but it does not
increase completed utilization: the additional bookings fill the finite
horizon, generate no-offer arrivals, and accumulate more daily cancellation
exposure. Thus, the regression coefficient should be read as a relative
loss-channel result, not as a simple claim that more balking is better.

\paragraph{Threshold ordering.}
Threshold ordering creates a separate mechanism from step intensity. The
average threshold level carries positive coefficients for utilization
\((\beta=+0.296)\), overall served rate \((\beta=+0.139)\), and mean offered
wait \((\beta=+0.148)\). Pushing both thresholds outward means that patients
can tolerate longer booking delays before entering the high-risk balking or
no-show regimes.

The threshold difference is more structural. The difference
\(\theta^{\text{balk}}-\theta^{\text{no-show}}\) has a negative coefficient
for utilization \((\beta=-0.259)\) and overall served rate \((\beta=-0.158)\).
At baseline,
\[
\theta^{\text{no-show}} < \theta^{\text{balk}},
\]
which creates an accept-then-miss region:
\[
\theta^{\text{no-show}} < \tau \leq \theta^{\text{balk}}.
\]
Patients in this interval still accept appointments, but they already face
elevated no-show risk. These bookings reserve capacity and may later be lost
on the service day. When the ordering is reversed, patients reach the balking
region before the no-show region, so fewer high-risk appointments are booked
and then missed. This is why threshold ordering matters beyond the absolute
level of the thresholds.

The baseline ordering therefore helps explain booked-slot losses, but it does
not imply that eliminating balking improves completed utilization. The direct
no-balking diagnostic shows that, under the current baseline, removing balking
mainly shifts losses toward no-offer and cancellations rather than increasing
completed visits.

\subsection{Class Spillovers}

\paragraph{Class gap drivers.}
The largest class-gap drivers are the cancellation probability gap
\((|\beta|=0.452)\), balking threshold gap \((|\beta|=0.429)\), and balking
step gap \((|\beta|=0.349)\). This is consistent with the pooled FCFS
formulation. At baseline, the two classes are symmetric, so class labels are
exchangeable up to simulation noise. Gaps appear when one class is made more
likely to accept long-delay offers, reject offers before booking, or cancel
after booking. Because both classes share the same booking pool, changing one
class's behavior also creates spillovers for the other class.

\mainfigure{regression_significant_standardized_coefficients.png}{%
Significant standardized regression coefficients from \(240\) randomized
parameter settings. Features and targets are standardized; error bars are
\(95\%\) HC3 robust confidence intervals. The dominance of total arrival rate
for overall served rate and mean offered wait reflects fixed-capacity
pressure, while class-gap coefficients are driven mainly by asymmetric
behavioral parameters. Full table in Appendix~\ref{app:regression}.
}{fig:regression}

%% ============================================================
\section{Balking Deep Dive: Selection vs.\ Access}
\label{sec:balk-deepdive}
%% ============================================================

This section isolates Class~1 balking while Class~2 stays at baseline, to
test whether lower accepted wait can accompany lower access.

\mainfigure{../../../../../outputs/class1_balking/figures/percent_serviced_by_class.png}{%
As Class~1 high-balking probability increases from 0 to~1, Class~1
served rate falls while Class~2 served rate rises. The congestion
relief that Class~2 gains is mechanically consistent with fewer
Class~1 patients booking long-delay slots.}{fig:balk-served}

\mainfigure{../../../../../outputs/class1_balking/figures/mean_accepted_delay_by_class.png}{%
Class~1 accepted wait falls as Class~1 high-balking probability
increases. This is a selection artifact: long-delay Class~1 patients
leave the accepted pool. Accepted wait should not be read as a
delay-improvement signal in this context.}{fig:balk-accepted}

The two figures together demonstrate a pattern the formulation produces
cleanly: accepted wait and served rate move in opposite directions
when balking changes. Any policy that raises balking step would show
apparently improved wait times while actually reducing the share of
Class~1 patients who complete service.

%% ============================================================
\section{Class-Gap Interpretation Under Pooled FCFS}
\label{sec:classgap}
%% ============================================================

Arrivals from both classes are pooled and randomly permuted each day before
FCFS assignment. The engine does not distinguish between classes when
offering slots. Therefore:

\begin{itemize}[leftmargin=1.5em, itemsep=3pt]
\item When class parameters are symmetric, the expected class served-rate
gap is zero. Observed deviations are simulation noise.
\item Meaningful class gaps arise only from asymmetric behavioral parameters
(different balking threshold, balking step, cancellation probability, or
no-show behavior) or from different arrival rates. These create systematic
differences in how often each class successfully books, cancels, and
attends.
\item A positive class~1 balking threshold relative to class~2 means
class~1 tolerates longer delays before rejecting offers. Under FCFS, a
more tolerant class captures more of the later-horizon slots and therefore
achieves a higher served rate.
\end{itemize}

\mainfigure{metric_class_gap_drivers.png}{%
When Class~1 behavioral parameters are varied while Class~2 stays at
baseline, a gap opens. Higher Class~1 cancellation, higher balking step,
or higher no-show step widen the gap against Class~1. Higher Class~1
balking threshold helps Class~1 by allowing it to accept longer-delay
slots. With symmetric parameters, both class lines coincide.
}{fig:class-gap}

The regression screen confirms this ranking: among class-gap drivers,
cancellation probability gap dominates ($|\beta| = 0.452$), followed by
balking threshold gap ($|\beta| = 0.429$) and balking step gap
($|\beta| = 0.349$). Class~1 arrival share enters significantly
($\beta = -0.120$): a larger Class~1 share tends to mildly reduce the
Class~1 served-rate advantage, consistent with more Class~1 arrivals
competing for the same slots under FCFS.

%% ============================================================
\section{What Seems Mechanically Robust vs.\ What Needs Validation}
\label{sec:checks}
%% ============================================================

\subsection{Mechanically Robust Under the Current Formulation}

\begin{itemize}[leftmargin=1.5em, itemsep=3pt]
\item \textbf{Utilization and served rate have different denominators.}
Their decoupling is a direct consequence of the metric definitions. It
is not a fragile numerical result.

\item \textbf{High utilization can coexist with low served rate.}
When demand exceeds supply, served rate is bounded by the supply-to-demand
ratio. This bound holds regardless of behavioral parameters.

\item \textbf{Suppressing balking mainly changes loss accounting.}
Under the baseline, suppressing balking does not materially change completed
utilization. It mostly changes where arrivals are lost: balked arrivals become
a combination of additional bookings, no-offer arrivals, and cancellations.

\item \textbf{Offered wait is less selective than accepted wait.}
Offered delay is recorded before balking, so offered wait includes patients
who ultimately leave. Accepted wait excludes them. The direction of the
difference is structural.

\item \textbf{No-offer is a separate access failure not captured by
offered wait.} No-offer patients receive no delay information and are
excluded from all delay metrics. Above $\lambda \approx 110$, they become
a substantial fraction of arrivals.

\item \textbf{Class labels are exchangeable under symmetric parameters.}
FCFS pooling and random permutation mean that class identity has no
systematic effect when classes are identical. The near-zero baseline class
gap (0.001) is consistent with this.

\item \textbf{The cancellation--delay feedback is directionally correct.}
Longer average booking delay creates more cancellation-check opportunities,
which raises cumulative cancellation probability. The high observed
cancellation fraction at long delays is mechanically consistent with this.

\item \textbf{No-offer is zero at baseline; access failures are behavioral.}
The booking horizon does not saturate at $\lambda = 100$ because
cancellations continuously free future capacity. Balking and cancellation,
not horizon saturation, are the primary failure modes.
\end{itemize}

\subsection{Needs Validation or Further Code Checks}

\vspace{6pt}
\begin{checkenv}
\begin{enumerate}[leftmargin=1.5em, itemsep=4pt]
\item \textbf{Validate the daily cancellation hazard assumption.}
Cancellation is implemented as a daily hazard over future bookings
(\texttt{simulation/engine.py}, lines~240--253): every future booking draws
Bernoulli(\texttt{cancel\_prob}) on each day before service. This makes
long-delay bookings much more likely to cancel before service and is central
to the no-balking diagnostic.

\item \textbf{Validate the no-show timing assumption.}
No-show is evaluated using the original booking delay $\tau$
(\texttt{booking.booking\_delay}; \texttt{simulation/engine.py},
lines~202--206), not residual waiting time. This makes the original offered
delay the no-show risk input at service time.

\item \textbf{Clarify no-offer propagation within a day.}
The documentation states that when the horizon is full for one patient,
all remaining arrivals that day are counted as \texttt{no\_offer}. This
means a single saturation event affects the full day remainder, which
could create day-order artifacts. Verify whether the no-offer share is
sensitive to the random permutation seed.

\item \textbf{Assess whether no-show slots should remain unbooked.}
No-show slots are not rebooked in the current formulation. In many
real clinics, same-day no-shows are filled with walk-ins or advance
overbooking. If this assumption is changed, both utilization and served
rate would shift; the magnitude depends on the no-show fraction and
same-day availability.

\item \textbf{Assess whether same-day cancellations should be excluded.}
Currently, no same-day cancellations are allowed. This keeps service-day
slots deterministic once the service day is reached. Relaxing this would
add a slot-recapture mechanism that could affect both utilization and
served rate.

\item \textbf{Confirm stability of class-gap estimates.}
The class-gap metric has the lowest out-of-sample $R^2$ (0.708) in the
regression screen and is near zero (0.001) at symmetric baseline. With
only 2~seeds per setting in the regression screen, sampling variance in
the gap could be material. At least 10--20 seeds per grid point are
advisable before drawing strong conclusions about class equity.

\item \textbf{Evaluate realism of baseline demand and capacity levels.}
100 arrivals per day against 32 slots per day is an intentional
stress-test, not a calibrated clinic. Before drawing any operational
conclusions, behavioral parameters (balking thresholds, no-show
probabilities, cancellation rates) should be estimated from clinic data
and the demand-to-capacity ratio brought into a plausible range.
\end{enumerate}
\end{checkenv}

%% ============================================================
\section{Summary for Discussion}
\label{sec:summary}
%% ============================================================

\begin{enumerate}[leftmargin=1.5em, itemsep=4pt]
\item Under the current formulation, utilization and served rate are
structurally decoupled. At baseline, this produces a 3:1 discrepancy
(0.839 vs.\ 0.269). Using utilization as a proxy for access would be
misleading at this demand level.

\item The dominant loss channels at baseline are balking (35.6\%) and
cancellation (32.3\%). No-offer is zero. Above $\lambda \approx 110$,
horizon saturation becomes an additional channel.

\item The controlled no-balking diagnostic shows that suppressing balking
increases booking conversion but not completed visits. Completed utilization
remains near~0.84; unmet demand shifts from balking to no-offer and
cancellations.

\item The high cancellation fraction is consistent with the long average
booking delay and the daily-cancellation rule. Whether this is realistic
depends on how real patients behave when booked many days out.

\item Offered wait is preferable to accepted wait as a delay measure,
because it includes patients who rejected the offer. But it still excludes
no-offer patients, who are the most severely access-constrained.

\item Class disparities in the simulation arise from asymmetric behavioral
parameters, not from FCFS favoritism. Symmetric classes produce
near-zero gaps.

\item The regression screen and one-at-a-time sensitivity runs are
internally consistent: demand pressure drives served rate and offered
wait; no-show and cancellation behavior drive utilization; behavioral
parameter gaps drive class equity.

\item None of these findings should be interpreted as causal estimates
about real appointment systems. They describe how this particular model
formulation responds to parameter variation.
\end{enumerate}

\bigskip
\noindent\textbf{Questions for discussion:}
\begin{enumerate}[leftmargin=1.5em, itemsep=4pt]
\item Is the baseline demand/capacity ratio intended as a stress scenario,
or should it be recalibrated toward a more plausible clinic setting?
\item Should cancellation remain a daily hazard over future bookings, or
should it be modeled as a one-time event at booking? The cumulative
cancellation fraction depends critically on this.
\item Should no-offer patients be reported as a separate access failure
in all summary tables, not only when the no-offer rate is non-negligible?
\item Should no-show probability depend on the original offered delay $\tau$
(current formulation) or on the remaining wait at service time?
\item Should same-day cancellations be allowed? What would the rebooking
mechanism look like?
\item Should alternative scheduling rules (e.g., priority-based,
overbooking) be compared against FCFS within this formulation?
\end{enumerate}

%% ============================================================
\appendix

\section{Exact Metric Definitions and Notation}
\label{app:defs}
%% ============================================================

Let $i$ be patient class, $D$ measured service days, $S$ slots per day,
$A_i$ arrivals, $B_i$ accepted bookings, $K_i$ balked patients, $C_i$
cancellations, $H_i$ no-shows, $V_i$ completed visits, and
$O_i = B_i + K_i$ offered patients.

\begin{align*}
\texttt{percent\_serviced}_i &= V_i / A_i, \\
\texttt{overall\_percent\_serviced} &= \textstyle\sum_i V_i \;/\; \textstyle\sum_i A_i, \\
\texttt{mean\_offered\_booking\_delay} &= \textstyle\sum_i \tau^{\text{off}}_i \;/\; \textstyle\sum_i O_i, \\
\texttt{mean\_accepted\_booking\_delay} &= \textstyle\sum_i \tau^{\text{acc}}_i \;/\; \textstyle\sum_i B_i, \\
\texttt{average\_utilization} &= \tfrac{1}{D}\textstyle\sum_d V_d / S, \\
\texttt{balking\_rate} &= \textstyle\sum_i K_i \;/\; \textstyle\sum_i O_i, \\
\Delta_{\text{access}} &= \texttt{percent\_serviced}_1 - \texttt{percent\_serviced}_2.
\end{align*}

\noindent\textbf{Note on terminology.} In the main text,
\texttt{percent\_serviced} (the simulation's internal variable name) is
reported as \emph{overall served rate}. \texttt{mean\_offered\_booking\_delay}
is reported as \emph{mean offered wait}; \texttt{mean\_accepted\_booking\_delay}
as \emph{mean accepted wait}.

\medskip
\noindent\textbf{Behavioral step functions.}
\[
  b_i(\tau) = \begin{cases}
    b_i^{\text{low}} & \tau < \theta_i^{\text{balk}}, \\
    b_i^{\text{high}} & \tau \geq \theta_i^{\text{balk}}.
  \end{cases}
  \qquad
  \xi_i(\tau) = \begin{cases}
    \xi_i^{\text{low}} & \tau < \theta_i^{\text{no-show}}, \\
    \xi_i^{\text{high}} & \tau \geq \theta_i^{\text{no-show}}.
  \end{cases}
\]

Balking step $= b_i^{\text{high}} - b_i^{\text{low}}$; no-show step
$= \xi_i^{\text{high}} - \xi_i^{\text{low}}$.

\medskip
\noindent\textbf{Measurement semantics.} Class metrics cover patients
arriving in the measurement window. If \texttt{cooldown\_days}~$<$~%
\texttt{horizon\_days}~$-$~1, some late bookings are unresolved at
simulation end and counted as \texttt{unresolved\_booked} (0.06\% of
arrivals at baseline). \texttt{total\_value} is service-day based and
is not cohort-aligned with class metrics.

%% ============================================================
\section{Simulation Design and Sensitivity Grid}
\label{app:sens-design}
%% ============================================================

All sensitivity grids are centered on \texttt{configs/baseline.yaml}.
One-axis slices vary Class~1 parameters while Class~2 stays at baseline.
Two-class heatmaps vary both classes independently.

\begin{table}[H]
\centering
\renewcommand{\arraystretch}{1.15}
\begin{tabular}{p{0.30\textwidth}rrr}
\toprule
Parameter & Baseline & Range & Grid size \\
\midrule
Arrivals (per class), 1D slices & 50 & 25--75 & continuous \\
Total arrivals, capacity stress & 100 & 30--170 & continuous \\
Balking step & 0.50 & 0.00--1.00 & 21 \\
Balking threshold (days) & 9 & 0--13 & 14 \\
No-show step & 0.30 & 0.00--1.00 & 21 \\
No-show threshold (days) & 6 & 0--13 & 14 \\
Cancellation probability & 0.10 & 0.00--0.30 & 21 \\
\bottomrule
\end{tabular}
\caption{Sensitivity parameter ranges. Each grid point uses seed 2027;
multi-seed comparisons average seeds 2027, 2028, 2029.}
\label{tab:sens-grid}
\end{table}

\textbf{Regression screen.}
240 randomized parameter settings (seed 9137); two seeds (4101, 4102)
averaged per setting; 12 features (total arrival rate, Class~1 share,
average level and Class~1$-$Class~2 gap for each behavioral parameter);
OLS with HC3 robust standard errors; only $p < 0.05$ coefficients shown.

\textbf{Balking deep dives.}
Class~1 high-balking-probability sweep over $\{0.0, 0.1, \ldots, 0.9\}$
(100 seeds per point; Class~2 fixed at baseline). Class~1 threshold sweep
over $\{3, 4, \ldots, 13\}$ (100 seeds per point; Class~2 fixed at
baseline).

%% ============================================================
\section{Supplementary Sensitivity Plots}
\label{app:sens}
%% ============================================================

\subsection{Driver Plots}

\mainfigure{metric_utilization_drivers.png}{%
Utilization across Class~1 arrival and behavioral parameter ranges.
Utilization stays near 0.839 across most of the demand range, while
class-level served rates diverge.}{fig:util-drivers}

\mainfigure{metric_access_drivers.png}{%
Served rate falls as arrival rates rise. The Class~1 and Class~2 lines
separate when Class~1 parameters diverge from symmetric baseline.
}{fig:access-drivers}

\mainfigure{metric_wait_drivers.png}{%
Offered and accepted wait across Class~1 parameter ranges. Accepted
wait can fall while offered wait stays elevated, illustrating selection
bias in the accepted-wait metric.}{fig:wait-drivers}

\mainfigure{metric_balking_rate_drivers.png}{%
Balking rate is controlled primarily by balking step and threshold.
Rising balking rate indicates congestion but not improved service.
}{fig:balk-drivers}

\subsection{No-Show and Cancellation Heatmaps}

\mainfigure{no_show_step_interaction_heatmap.png}{%
No-show step heatmap. Utilization falls when both classes carry high
no-show risk. Class gap moves against the higher no-show class.
}{fig:noshow-step}

\mainfigure{no_show_threshold_interaction_heatmap.png}{%
No-show threshold heatmap. Higher threshold delays high no-show onset,
improving utilization and access for the affected class.
}{fig:noshow-threshold}

\mainfigure{cancellation_probability_interaction_heatmap.png}{%
Cancellation heatmap. Higher cancellation for one class reduces that
class's access; aggregate effects are more complex (see
Section~\ref{sec:sensitivity}).}{fig:cancel}

\subsection{Balking and No-Show Interaction Heatmaps}
\label{app:balk-noshow-heatmaps}

\mainfigure{balk_noshow_step_utilization.png}{%
Average utilization as a function of balking step and no-show step, with
baseline slices shown in the side panels. The nearly flat balking-step slice
at fixed no-show step supports the interpretation that no-show exposure is
the more direct utilization channel.
}{fig:balk-noshow-heatmap}

\mainfigure{balk_noshow_threshold_utilization.png}{%
Average utilization as a function of balking threshold and no-show threshold.
The baseline lies in the region where
$\theta^{\text{no-show}} < \theta^{\text{balk}}$, creating an
accept-then-miss interval for long-delay bookings.
}{fig:balk-noshow-threshold-heatmap}

\subsection{Balking Heatmaps}

\mainfigure{balking_step_interaction_heatmap.png}{%
Balking step heatmap. Served rate falls and class gap widens as one
class's balking step rises above the other.}{fig:balk-step}

\mainfigure{balking_threshold_interaction_heatmap.png}{%
Balking threshold heatmap. The more tolerant class (higher threshold)
captures more served slots under FCFS.}{fig:balk-threshold}

\mainfigure{arrival_mix_interaction_heatmap.png}{%
Arrival mix heatmap. Left: overall served rate. Right: mean offered wait.}{fig:arrival-mix}

\subsection{Capacity Stress}

\mainfigure{fcfs_capacity_stress_curves.png}{%
Capacity stress diagnostic curves. Served rate and offered wait as
demand rises from $\lambda = 30$ to~170; utilization remains nearly
flat across the full range.}{fig:stress}

\subsection{Balking Deep Dive: Additional Plots}

\mainfigure{../../../../../outputs/class1_balking/figures/mean_offered_delay_by_class.png}{%
Offered wait by class as Class~1 high-balking probability rises.
When enough Class~1 patients reject long-delay offers, overall horizon
congestion falls and both classes receive shorter offered delays.
}{fig:balk-offered}

\mainfigure{../../../../../outputs/class1_balking_threshold/figures/percent_serviced_by_class.png}{%
Class~1 balking threshold sweep: higher threshold increases Class~1
access by tolerating longer delays. The gain comes at mild cost to
Class~2 when Class~1 tolerance is very high.}{fig:thresh-served}

\mainfigure{../../../../../outputs/class1_balking_threshold/figures/service_gap.png}{%
Class~1 minus Class~2 served-rate gap as Class~1 threshold varies.
}{fig:thresh-gap}

%% ============================================================
\section{Regression Screen}
\label{app:regression}
%% ============================================================

\begin{table}[H]
\centering
\small
\renewcommand{\arraystretch}{1.08}
\begin{tabular}{p{0.17\textwidth}p{0.27\textwidth}rp{0.20\textwidth}r}
\toprule
Target & Feature & Std.\ coef. & 95\% HC3 CI & HC3 $p$ \\
\midrule
Utilization & Avg threshold level (balk $+$ no-show) & $+0.296$ & $[0.219,\;0.373]$ & $<0.001$ \\
Utilization & Balk step $-$ no-show step & $+0.365$ & $[0.227,\;0.504]$ & $<0.001$ \\
Utilization & Avg cancellation prob. & $+0.311$ & $[0.240,\;0.384]$ & $<0.001$ \\
Utilization & Avg step level (balk $+$ no-show) & $-0.289$ & $[-0.371,\;-0.207]$ & $<0.001$ \\
Utilization & Balk threshold $-$ no-show threshold & $-0.259$ & $[-0.387,\;-0.131]$ & $<0.001$ \\
Utilization & Sign(balk threshold $-$ no-show threshold) & $-0.198$ & $[-0.330,\;-0.067]$ & $0.003$ \\
Utilization & Total arrival rate & $-0.104$ & $[-0.182,\;-0.027]$ & $0.008$ \\
\midrule
Overall served rate & Total arrival rate & $-0.773$ & $[-0.854,\;-0.692]$ & $<0.001$ \\
Overall served rate & Avg threshold level (balk $+$ no-show) & $+0.139$ & $[0.084,\;0.195]$ & $<0.001$ \\
Overall served rate & Balk step $-$ no-show step & $+0.174$ & $[0.077,\;0.271]$ & $<0.001$ \\
Overall served rate & Avg step level (balk $+$ no-show) & $-0.107$ & $[-0.168,\;-0.045]$ & $0.001$ \\
Overall served rate & Balk threshold $-$ no-show threshold & $-0.158$ & $[-0.264,\;-0.053]$ & $0.003$ \\
Overall served rate & Avg cancellation prob. & $+0.114$ & $[0.054,\;0.174]$ & $<0.001$ \\
\midrule
Mean offered wait & Total arrival rate & $+0.575$ & $[0.496,\;0.653]$ & $<0.001$ \\
Mean offered wait & Avg cancellation prob. & $-0.440$ & $[-0.504,\;-0.375]$ & $<0.001$ \\
Mean offered wait & Balk threshold $-$ no-show threshold & $+0.173$ & $[0.049,\;0.296]$ & $0.006$ \\
Mean offered wait & Avg threshold level (balk $+$ no-show) & $+0.148$ & $[0.069,\;0.228]$ & $<0.001$ \\
Mean offered wait & Avg step level (balk $+$ no-show) & $-0.181$ & $[-0.271,\;-0.091]$ & $<0.001$ \\
\midrule
Class gap & Cancellation prob.\ gap & $-0.452$ & $[-0.558,\;-0.346]$ & $<0.001$ \\
Class gap & Balking threshold gap & $+0.429$ & $[0.331,\;0.528]$ & $<0.001$ \\
Class gap & Balking step gap & $-0.352$ & $[-0.436,\;-0.268]$ & $<0.001$ \\
Class gap & No-show threshold gap & $+0.266$ & $[0.175,\;0.357]$ & $<0.001$ \\
Class gap & No-show step gap & $-0.213$ & $[-0.289,\;-0.136]$ & $<0.001$ \\
Class gap & Class~1 arrival share & $-0.119$ & $[-0.209,\;-0.029]$ & $0.009$ \\
\bottomrule
\end{tabular}
\caption{Significant standardized regression coefficients ($p<0.05$, HC3 robust SEs).
Gap = Class~1 minus Class~2 value. 240 randomized settings, 2 seeds averaged per
setting. Only features passing $p < 0.05$ are shown.}
\label{tab:regression}
\end{table}

\mainfigure{regression_significant_standardized_coefficients.png}{%
Significant standardized regression coefficients with 95\% HC3 confidence
intervals. Total arrival rate dominates served rate and offered wait;
no-show behavior dominates utilization; behavioral parameter gaps drive
the class served-rate gap.}{fig:app-regression}

%% ============================================================
%% Changelog
%% ============================================================
% Base file: metric_analysis_research_note.tex (v1)
% Sections rewritten:
%   - §2 Baseline (added formal parameter table)
%   - §10 Summary (renamed "Summary for Discussion"; added discussion questions)
% Sections added:
%   - §4 Cross-Metric Diagnostics (from merged)
%   - App B Sensitivity Design (from merged)
%   - App D Regression Screen (from merged)
% Figures in main body:
%   - fcfs_arrival_outcome_decomposition.png (§3)
%   - research_note_main_panel.png (§4)
%   - regression_significant_standardized_coefficients.png (§6)
%   - class1_balking deep-dive figures (§7)
%   - metric_class_gap_drivers.png (§8)
% Figures moved to appendix:
%   - metric_utilization_drivers.png, metric_access_drivers.png,
%     metric_wait_drivers.png, metric_balking_rate_drivers.png (App C)
%   - all heatmaps (App C)
%   - fcfs_capacity_stress_curves.png (App C)
%   - additional deep-dive threshold figures (App C)
% Figures removed:
%   - scenario_metric_comparison.png (Scenario 2 reference)
% Scenario 2: Removed entirely.
%   Reason: configs/scenario_2.yaml changes 6 parameters simultaneously
%   (horizon, lambda_2, balk_threshold, balk_high, no_show_threshold,
%   no_show_high). Not a one-parameter comparison. util=1.000 in merged
%   table is suspicious. Per hard constraint #2, removed.
% Unresolved validation checks (carried as checkenv callouts):
%   1. Cancellation: once per residual day vs. once per booking
%   2. No-offer propagation creates daily-order artifacts
%   3. Positive cancellation coefficient on served rate
%   4. Class-gap stability with more seeds

\end{document}
